% Version 0.24.  Classical finite duality, with all source conventions matched.
\section{Finite duality and unramified reciprocity on the original fields}
\label{sec:local-duality-closure}

\status{This section supplies the literal finite-module
\texttt{UnramifiedOrthogonality} interface and matches
\texttt{Published.PoitouTateDuality} to the classical finite-module
Poitou--Tate theorem.  The elementary procyclic, Euler-characteristic,
annihilator and reciprocity deductions are proved.  Local Tate duality,
local class field theory and finite-module Poitou--Tate duality are
explicit classical citation leaves.  No main conjecture, special-value
estimate, or arithmetic prime window is assumed or proved here.}

The relevant coefficients are finite modules, not necessarily free
modules.  Fix a prime $\ell$, an integer $r>0$, and
$R=\mathbb Z/\ell^r\mathbb Z$.  For a number field $F$ with its specified
algebraic closure $\Omega$, write $R(1)=\mu_{\ell^r}(\Omega)$, in additive
notation.  A primitive root $\zeta$ gives the identification
$a\mapsto\zeta^a$ with the source's rank-one representation
\texttt{charRep(cycMod)}.  It is equivariant because
$g\zeta=\zeta^{\chi_{\ell,r}(g)}$.  The character
$\chi_{\ell,r}$ is independent of the chosen primitive root.
For a finite continuous $R$-representation $M$, put
\[
 M^D=\operatorname{Hom}_R(M,R(1)),\qquad
 (g\varphi)(m)=\chi_{\ell,r}(g)\varphi(g^{-1}m).
 \tag{LD.1}
\]
Thus $M^D$ is exactly \texttt{twDual}, including its inverse in the
argument of $\varphi$.

At a finite place $v$, keep the chosen valuation subring $A\subset\Omega$.
Its decomposition group $D_A$ and inertia group $I_A$ identify with
$G_{F_v}$ and its inertia.  One can make this identification by extending
the chosen valuation to an embedding in an algebraic closure of $F_v$:
finite extensions of the henselization and of the completion have the
same finite separable extensions, with the same ramification and residue
extensions.  Passing to inverse limits gives the asserted group
identification.  A different embedding conjugates both groups and every
restriction map.  We use the given $A$ throughout; this identification
does not replace $F$, $\Omega$, or its local data by a new arithmetic
witness.

\subsection{Procyclic cohomology and the Euler calculation}

\begin{proposition}[LD1: finite procyclic and tame Euler calculation]
\label{prop:ld-procyclic}
Let $\Gamma=\widehat{\mathbb Z}$ with its specified topological generator
$\phi$, acting continuously on a finite $\ell$-primary module $T$.  Then
\[
 H^0(\Gamma,T)=\ker(\phi-1),\qquad
 H^1(\Gamma,T)=T/(\phi-1)T,\qquad
 H^i(\Gamma,T)=0\quad(i\ge2).
 \tag{LD.2}
\]
In particular $|H^0(\Gamma,T)|=|H^1(\Gamma,T)|$.
If $K$ is a nonarchimedean local field of characteristic zero and
residue characteristic $p\ne\ell$, and $M$ is any finite continuous
$\ell$-primary $G_K$-module, then its cohomology is finite and
\[
 |H^1(K,M)|=|H^0(K,M)|\,|H^2(K,M)|,
 \qquad H^i(K,M)=0\quad(i\ge3).
 \tag{LD.3}
\]
No unramifiedness assumption on $M$ is needed in (LD.3).
\end{proposition}

\begin{proof}
A continuous crossed homomorphism on $\Gamma$ is determined by its
value at $\phi$, since the integer powers of $\phi$ are dense.  Every
$t\in T$ occurs: if $\phi^d$ acts trivially and $\ell^aT=0$, the rule
\[
 f(\phi^j)=\sum_{i=0}^{j-1}\phi^it\quad(j\ge0)
\]
factors through the cyclic quotient of order $d\ell^a$.  Indeed its
value at that order is $\ell^a\sum_{i=0}^{d-1}\phi^it=0$.
Coboundaries take at $\phi$ the values $(\phi-1)t$, proving the
degree-one formula.  Invariants give degree zero.

For completeness, the higher-degree calculation is the continuous
version of the cyclic resolution with alternating operators
$\phi-1$ and $1+\phi+\cdots+\phi^{n-1}$.  A continuous cochain with
finite target factors through a finite cyclic quotient; hence continuous
cohomology is the direct limit over such quotients.  Inflation from
order $n$ to order $kn$, where both actions factor through order $n$,
multiplies the even-degree norm-class representative by $k$ in degree
two.  In every degree $i\ge2$ the comparison of the periodic
resolutions contains a factor $k$ (and possibly a higher power of $k$).
Choosing $k$ divisible by the exponent of $T$ kills that class.
This proves the vanishing in (LD.2).  Alternatively the same calculation
is the two-term continuous resolution for the pro-$\ell$ factor,
after taking invariants under the prime-to-$\ell$ factor.  Finally the
kernel and cokernel of an endomorphism of a finite group have equal
orders, which proves the cardinality assertion.

For (LD.3), let $I$ be inertia and $P$ wild inertia.  The group $P$ is
pro-$p$, so its higher cohomology on $\ell$-primary modules vanishes:
on each finite quotient the usual averaging homotopy divides by a
$p$-power, which is invertible on the module.  The tame quotient is
\[
 I/P\simeq\prod_{q\ne p}\mathbb Z_q(1).
\]
The prime-to-$\ell$ factor again has exact invariants.  Applying the
procyclic calculation to its remaining $\mathbb Z_\ell(1)$ factor
shows that $A=M^I$ and $B=H^1(I,M)$ are finite and $H^i(I,M)=0$ for
$i\ge2$.  The residue Galois group $G_K/I$ is
$\widehat{\mathbb Z}$.  The Hochschild--Serre cochain filtration,
with only its two inertia rows and two residue columns remaining, gives
\[
 H^0(K,M)=A^\phi,\qquad
 0\longrightarrow A/(\phi-1)A\longrightarrow H^1(K,M)
   \longrightarrow B^\phi\longrightarrow0,
 \tag{LD.4}
\]
and $H^2(K,M)=B/(\phi-1)B$, with no higher groups.  All possible
differentials leave these two columns, so are zero.  Applying the
kernel--cokernel order equality separately to $A$ and $B$ gives
(LD.3).  This also proves all required finiteness in this case.
\end{proof}

\subsection{The exact annihilator of unramified classes}

\begin{proposition}[LD2: the literal unramified orthogonality supplier]
\label{prop:ld-unramified}
Let $K=F_v$ have residue characteristic different from $\ell$, and
let $M$ be a finite continuous $R$-representation on which inertia acts
trivially.  Define
\[
 H^1_{\rm ur}(K,M)=
 \ker\bigl(H^1(K,M)\longrightarrow H^1(I,M)\bigr).
\]
In the finite local Tate pairing
\[
 H^1(K,M^D)\times H^1(K,M)
 \xrightarrow{\ \cup\ }H^2(K,R(1))\simeq R,
 \tag{LD.5}
\]
the two unramified subgroups are exact annihilators in both directions.
Consequently \texttt{ReciprocityTheory.Gal.UnramifiedOrthogonality}
holds for every original number field, algebraic closure, place and
finite coefficient module occurring in that interface.
\end{proposition}

\begin{proof}
Prime-to-$p$ roots of unity are unramified over $K$.  To see this,
choose a finite residue extension containing the roots of
$X^{\ell^r}-1$; its derivative is invertible at every root.  Hensel
lifting in the corresponding unramified local extension gives all
$\ell^r$-th roots.  Thus $R(1)$, and hence $M^D$, are unramified.

Inflation--restriction and LD1 identify
\[
 H^1_{\rm ur}(K,M)=H^1(G_K/I,M)=M/(\phi-1)M,
 \qquad |H^1_{\rm ur}(K,M)|=|H^0(K,M)|.
 \tag{LD.6}
\]
The same holds for $M^D$.  The product of two unramified classes is
inflated from $H^2(G_K/I,R(1))=0$.  Each unramified subgroup therefore
annihilates the other.

The one classical input in this argument is finite local Tate duality:
cup product followed by the local invariant is a perfect pairing
$H^i(K,M)\times H^{2-i}(K,M^D)\to\mathbb Q/\mathbb Z$,
for $i=0,1,2$.  See Milne, \emph{Arithmetic Duality Theorems},
second edition, I, Corollary~2.3.  Since these groups are killed by
$\ell^r$, the target is its $\ell^r$-torsion, identified with $R$.
In particular
\[
 |H^2(K,M)|=|H^0(K,M^D)|.
\]
Combine this with LD1 and (LD.6):
\[
 |H^1(K,M)|=
 |H^1_{\rm ur}(K,M)|\,|H^1_{\rm ur}(K,M^D)|.
 \tag{LD.7}
\]
For a perfect pairing of finite abelian groups $X$ and $Y$, the
annihilator in $X$ of $Y_0\subset Y$ has order $|X|/|Y_0|$.
Indeed every character of $Y_0$ extends to $Y$, since
$\mathbb Q/\mathbb Z$ is divisible, so the restriction map from
$Y^\vee$ onto $Y_0^\vee$ is surjective.  Equation (LD.7) makes the
already proved inclusion an equality.  Interchanging $M$ and $M^D$
proves the other equality.

There is no hidden freeness condition.  A finite $R$-module is a direct
sum of $R/(\ell^a)$ with $0<a\le r$.  Its $R$-dual on that summand is
$\ell^{r-a}R$, of the same order, and evaluation identifies it with
its double dual.  Thus the $R$-linear duality used in the source agrees
with the finite abelian-group duality above.  The source places its
dual class first in (LD.5); reversing degree-one factors only changes
a sign, which does not change a zero or an annihilator.  Finally the
source requires vanishing in $H^2(K,R(1))$, before applying an
invariant.  The invariant is an isomorphism, so this is exactly the
same condition.  All identifications were made on the originally
chosen decomposition and inertia groups.
\end{proof}

\subsection{One normalized family of local invariants}

\begin{proposition}[LD3: exact finite Poitou--Tate specialization]
\label{prop:ld-poitou-tate}
For every number field $F$, specified algebraic closure $\Omega$,
prime $\ell$, integer $r>0$, and primitive $\ell^r$-th root
$\zeta\in\Omega$, assume $\ell\ne2$ or $F$ is totally complex.
The classical finite-module Poitou--Tate theorem supplies the exact
interface \texttt{GlobalTheory.Published.PoitouTateDuality}.
In particular its local invariant family can be fixed \emph{before}
the finite set $S$, all choices of places over $S$, the finite module,
and its continuous action are chosen.
\end{proposition}

\begin{proof}
We first match the normalization, rather than choosing unrelated
isomorphisms $H^2\simeq R$ at different places.  For each finite place
of $F$, Kummer theory embeds $H^2(F_v,\mu_{\ell^r})$ in the Brauer
group with image its $\ell^r$-torsion.  The local invariant
$\operatorname{inv}^{\rm Br}_v:\operatorname{Br}(F_v)\simeq
\mathbb Q/\mathbb Z$ identifies this image with
$\ell^{-r}\mathbb Z/\mathbb Z$.  Let
$\iota_r:R\to\mathbb Q/\mathbb Z$ be
$a\mapsto a/\ell^r$.

Use the common sign for which the Kummer class of a uniformizer,
cupped on the left with the unramified character taking arithmetic
Frobenius to $1$, has invariant $1\in R$.  More explicitly, with the
arithmetic local reciprocity convention in which
\[
 \operatorname{inv}^{\rm Br}_v(\chi\cup\kappa(\varpi))
       =\iota_r(\chi(\operatorname{Frob}_v)),
\]
the source's order $\kappa(\varpi)\cup\chi$ requires
$j_v=-\iota_r^{-1}\operatorname{inv}^{\rm Br}_v$.
The sign follows from graded commutativity.  This elementary cyclic
algebra normalization is part of local class field theory; it is
independent of $v$, $S$, and $M$.

For a chosen valuation subring $A$ above $v$, transport this invariant
along its decomposition-group identification and the map
$a\mapsto\zeta^a$.  The formula is a bijective $R$-linear map
$j_A:H^2(D_A,R(1))\to R$.  Its uniformizer normalization is literally
the source's \texttt{IsNormalised}: a Kummer cocycle $\kappa$ and an
unramified additive character $\chi$ have the displayed representative
$\kappa(g)\chi(h)$, and arithmetic Frobenius is the residue power
map used in \texttt{IsFrob}.  At valuation subrings which do not lie
over finite places, extend the family by the zero map; the interface
places no condition there.

This family is also \texttt{ConjCompatible}.  If $\sigma$ transports
$A$ to $A'$ and $\sigma\zeta=\zeta^c$, the coefficient description
of a transported cocycle is $z'=cz$.  Under the root-of-unity
identification, this says precisely
$\zeta^{z'}=\sigma(\zeta^z)$.  The Brauer class is therefore
transported by the induced field isomorphism, which preserves the
local invariant.  Hence $j_{A'}([z'])=j_A([z])$, the exact source
identity, including the coefficient $c$.  All choices in the family
have now been made.

Let $S$ be an arbitrary finite set of finite places containing all
places above $\ell$.  Enlarge it by the archimedean places only for
the statement of the classical theorem, and write $G_{F,S}$ for the
Galois group unramified outside its finite places.  The source's
quotient by the closed normal subgroup generated by the omitted
inertia groups is precisely this group: a finite quotient factors
through it if and only if its finite extension is unramified at
every omitted finite place.  The coefficient $R(1)$ is unramified
there by the root-of-unity argument in LD2, and hence descends to
$G_{F,S}$.  If $\mu$ is any character whose pullback is the original
cyclotomic character, it is this same descended character, since
the quotient map is onto.  In particular the source's
\texttt{toMuLoc} is induced by the identity on $R$, not by an
uncontrolled new coefficient identification.

We now apply Milne, \emph{Arithmetic Duality Theorems}, second
edition, I, Theorem~4.10, to an arbitrary finite continuous
$R[G_{F,S}]$-module $M$.  The order of $M$ is a power of $\ell$,
so it is a unit in $\mathcal O_{F,S}$, as required by that theorem.
Let $P^i_S(M)=\prod_{v\in S}H^i(F_v,M)$, for finite $S$, and let
$T^\vee=\operatorname{Hom}_R(T,R)$.  The resulting nine-term sequence
is
\[
\begin{aligned}
0\longrightarrow H^0(G_{F,S},M)&\longrightarrow P^0_S(M)
 \longrightarrow H^2(G_{F,S},M^D)^\vee\\
 &\longrightarrow H^1(G_{F,S},M)
 \longrightarrow P^1_S(M)
 \longrightarrow H^1(G_{F,S},M^D)^\vee\\
 &\longrightarrow H^2(G_{F,S},M)
 \longrightarrow P^2_S(M)
 \longrightarrow H^0(G_{F,S},M^D)^\vee\longrightarrow0.
\end{aligned}
 \tag{LD.8}
\]
Here the maps from the local terms are the sums of the local
cup-product pairings followed by the \emph{same} $j_A$ just fixed.
The source's middle pairing reverses the degree-one factors, which
multiplies that map by $-1$ relative to the opposite convention.
Changing the common invariant sign also multiplies the transgression
maps by $-1$.  Such changes preserve their images and kernels, and
the connecting maps may be given the corresponding signs.  Thus
the exact sequence has precisely the source's
\texttt{loc0,loc1,loc2} and \texttt{tr0,tr1,tr2}; it asserts the
existence of its two $R$-linear connecting maps.  This is
\texttt{NineTermExact}, not merely equality of orders.

The classical local terms at infinity use Tate cohomology, including
degree zero.  At a complex place they vanish.  At a real place they
vanish on an odd-primary module because averaging over the group of
order two is exact and its norm is onto the invariants.  Under the
assumption in the proposition, all these terms therefore vanish
and (LD.8) has exactly the finite-place products in the source.
Ordinary $H^0$ at a real place must not be inserted here.

For $i=1,2$ let $\operatorname{Sha}^i_S(M)$ be the kernel of localization to the
finite places in $S$.  Part (a) of the same classical theorem gives
perfect pairings of finite groups
\[
 \operatorname{Sha}^1_S(M)\times\operatorname{Sha}^2_S(M^D)\longrightarrow\mathbb Q/\mathbb Z,
 \qquad
 \operatorname{Sha}^1_S(M^D)\times\operatorname{Sha}^2_S(M)\longrightarrow\mathbb Q/\mathbb Z.
 \tag{LD.9}
\]
They take values in $\ell^{-r}\mathbb Z/\mathbb Z$ and are
$R$-bilinear.  The finite dual-module calculation in LD2 converts
them into the two $R$-linear equivalences and equal cardinalities
specified by \texttt{ShaDual}.  Part (c) identifies higher global
cohomology with the real-place groups; those vanish by the same
argument, so $H^n(G_{F,S},M)=0$ for all $n\ge3$.
This verifies every clause and the order of every quantifier in
\texttt{PoitouTateDuality}.
\end{proof}

\subsection{The finite reciprocity consequence actually used}

\begin{proposition}[LD4: global cup reciprocity and one exceptional place]
\label{prop:ld-reciprocity}
Retain LD3's original field, finite set $S$ and local invariant
family.  For every $z\in H^2(G_{F,S},R(1))$,
\[
                    \sum_{v\in S}j_{A_v}(z|_{D_{A_v}})=0.
                    \tag{LD.10}
\]
Consequently, if all but one localization of $z$ vanish, the remaining
one vanishes too.  In particular this applies to the cup product of
two global finite cocycles with equivariant pairing into $R(1)$.
At a place away from $\ell$ where the coefficient representations
are unramified, two unramified local classes contribute zero; and
orthogonality to every unramified class in the Tate-dual module
forces a class itself to be unramified.
\end{proposition}

\begin{proof}
Apply (LD.8) to $M=R(1)$.  Its dual is the trivial rank-one module
$R$, and the identity functional is the global invariant $1$ in
$H^0(G_{F,S},R)$.  Exactness at $P^2_S(R(1))$ says that the image of
any global $z$ is killed by the last transgression.  Evaluating that
functional at $1$ gives exactly (LD.10).  The termwise description
uses evaluation, so no trace factor or degree appears.  If only
one term remains, its local invariant is zero and its class is
zero because $j_{A_v}$ is bijective.  Naturality of cup products
under restriction gives the assertion for global cocycles.  A local
coboundary contributes zero to cohomology; a pair of unramified
classes contributes zero by LD2.  The last assertion is the reverse
annihilator implication of LD2.

These arguments can be made at any finite coefficient level before
taking a limit.  A continuous action on a finite module and finitely
many continuous cocycles factor through a finite Galois quotient:
local constancy on the compact profinite domain supplies a common
open normal subgroup.  Only finitely many finite places ramify in
that extension.  Adding them and the places above $\ell$ to $S$
therefore puts all these objects in the same $G_{F,S}$.  Enlarging
$S$ does not change the already chosen invariants and merely adds
zero unramified contributions at the new places where appropriate.
This is the global/single-exception input used by the source's
\texttt{OverIff.glob} and its finite reciprocity construction.
\end{proof}

\paragraph{Classical references and the exact closed boundary.}
The primary reference is J.~S.~Milne,
\href{https://jmilne.org/math/Books/ADTnot.pdf}{\emph{Arithmetic Duality
Theorems}, second edition (2006)}: I, Corollary~2.3 for finite local
Tate duality, I, Theorem~2.6 for the classical unramified-annihilator
statement, I, Lemma~2.9 for the prime-to-residue-characteristic Euler
identity, and I, Theorem~4.10 for finite Poitou--Tate duality.
LD1--LD2 provide direct proofs of the latter two local deductions
from the stated local duality theorem.  The Kummer/unramified
cyclic-algebra invariant is the standard local class field
normalization; the order of its factors and its sign were fixed
explicitly above.  For this normalization see Milne,
\href{https://www.jmilne.org/math/CourseNotes/CFT.pdf}{\emph{Class Field
Theory}, version 4.03 (2020)}, III, Proposition~3.6, Lemma~3.7,
and the cup-product construction in III, Section~4, Steps~1--3.
The author's
\href{https://www.jmilne.org/math/Books/add/ADT2006.pdf}{published
errata} and the exact finite-module hypotheses were checked.
These citations are established general duality results, not
unexpanded bespoke special-value bounds.

The supplied source endpoints are \texttt{UnramifiedOrthogonality}
and \texttt{Published.PoitouTateDuality}, with their full original-field,
finite-module and local-choice quantifiers.  LD4 records their finite
reciprocity consequence.  The Jacobi-sum/Kolyvagin and generalized
Bernoulli argument needed for \texttt{PsiBound} is a separate
arithmetic step; none of (LD.2)--(LD.10) asserts that special-value
bound on its own.
